\begin{textsample}{POS Dim 3 – pa – Score 15.00 – pa/pa\_inf\_16.txt}  \label{ex:f3_pos_271}
1.
O EU, O OUTRO E O NÓS O primeiro campo de experiência define a importância de se perceber em sua individualidade, suas características, emoções e sensações, os \textbf{cuidados} pessoais.
Nele ainda é informado sobre o processo de interação com outras \textbf{crianças} e com os \textbf{adultos}, destacando a percepção da \textbf{criança} de se relacionar afetiva e respeitosamente com o outro, compreendendo as diferentes culturas, costumes e as diferentes funções cumpridas pelos sujeitos.
Na interligação com o Ensino Fundamental, o campo de experiência “ O eu, o outro e o nós ” está particularmente relacionado ao eixo estruturante “ Valores à Vida Social ”, pois nesse eixo se manifestam os valores que constituem a sociedade por meio das relações sociais, em outros termos, da interação entre os sujeitos. 2.
CORPO, \textbf{GESTOS} E MOVIMENTOS Este campo de experiência traduz o corpo como estrutura física não só da \textbf{criança} como ser vivo, uma estrutura que possui habilidades importantes e necessárias para o desenvolvimento da aprendizagem, pois é na corporeidade que a exploração do mundo e das coisas se efetiva, seja por meio “ dos sentidos, \textbf{gestos}, movimentos impulsivos ou intencionais, coordenados ou espontâneos ” ( BRASIL, 2017a, p. 38 ).
O pleno desenvolvimento da \textbf{criança} se legitima a partir dos \textbf{gestos} e movimentos que ela apresenta em sua \textbf{rotina}, pois o corpo é por excelência um instrumento de comunicação e emancipação da \textbf{criança}.
Elas conhecem e reconhecem as sensações e funções de seu corpo [... ] identificam suas potencialidades e seus limites, desenvolvendo, ao mesmo tempo, a consciência sobre o que é seguro e o que pode ser um risco à sua integridade física ( BRASIL, 2017a, p. 39 ).
Na interligação com o Ensino Fundamental, o campo de experiência “ Corpo, \textbf{gestos} e movimentos ” está diretamente relacionado ao eixo estruturante “ Linguagem e suas Formas Comunicativas”[^21 ], visto que nesse eixo a linguagem é fator essencial para o desenvolvimento humano na forma de comunicação, pois serve para expor sentimentos, emoções e informações verbais, corporais, artísticas e dos sonidos. 3.
TRAÇOS, \textbf{SONS}, CORES E FORMAS Este campo de experiência discorre da importância de promover o convívio da \textbf{criança} com diversas manifestações artísticas, culturais e científicas ; regionais ou globais, objetivando explorar o senso estético, pela sensibilidade e \textbf{curiosidade} da \textbf{criança}.
Destaca -se também pela pertinência em ampliar o repertório cultural da \textbf{criança} diversificando o conhecimento acerca das culturas existentes, seja indígena, quilombola, ribeirinha, rural, urbana, africanas, europeias, asiáticas ou americanas.
Nesse campo se indica que as interações com as culturas citadas podem ser expressas para e pela \textbf{criança} por meio de “ diversas formas de expressão e linguagens, como as artes visuais ( \textbf{pintura}, modelagem, colagem, fotografia etc. ), a música, o teatro, a dança e o audiovisual, entre outras ”. ( BRASIL, 2017a, p. 39 ).
Na interligação com o Ensino Fundamental, o campo de experiência “ Traços, \textbf{sons}, cores e formas ” está relacionado com o eixo estruturante “ Cultura e Identidade ”, visto que esse eixo apresenta como indicativo o entendimento e interpretação das identidades que compõem as diversas culturas, ressaltando as relações sociais, dos sujeitos com o mundo. 4. \textbf{ESCUTA}, FALA, PENSAMENTO E IMAGINAÇÃO Este campo de experiência menciona que as interações vivenciadas pela \textbf{criança} desde que são \textbf{bebês} possibilitam situações comunicativas presentes em seus cotidianos que se apresentam nos movimentos de seus corpos, choro, \textbf{balbucio}, sorrisos, gargalhadas, olhares e que, com o crescimento delas, possibilitam ampliar e melhor desenvolver suas habilidades de comunicação.
Menciona ainda que atenção e \textbf{curiosidade} com a cultura escrita também se apresentam na vida dos \textbf{pequenos} desde que são \textbf{bebês}, devendo ser estimuladas pelas instituições de Educação \textbf{Infantil} para que a \textbf{criança} sinta prazer e familiaridade com o mundo da leitura e da escrita.
Esse campo também informa que a \textbf{escuta}, a fala, o pensamento e a imaginação são aguçados quando a \textbf{criança} mantém contato regular com a literatura \textbf{infantil}, as histórias, cordéis, músicas, poemas, fábulas que contribuem para o seu desenvolvimento afetivo, social e cognitivo.
Na interligação com o Ensino Fundamental, o campo de experiência “ \textbf{Escuta}, fala, pensamento e imaginação ”, assim como o campo “ Corpo, \textbf{gestos} e movimentos ”, está diretamente relacionado ao eixo estruturante “ Linguagem e suas Formas Comunicativas ”, pois escutar, falar, pensar e imaginar são ações que se justificam pela interação e comunicação entre os sujeitos. 5.
ESPAÇOS, TEMPOS, QUANTIDADES, RELAÇÕES E TRANSFORMAÇÕES O último campo de experiência entende que, por serem sujeitos históricos que se constituem em universos variados, essa posição histórica e geográfica contribui para que a \textbf{criança} se perceba dentro de espaços diversos ( rua, bairro, cidade, país, estado ) e diferentes tempos ( dia e noite, ontem, hoje e amanhã ) ; e ainda, por apresentar a \textbf{curiosidade} como característica peculiar de \textbf{criança}, ela tende a querer entender esse mundo da qual faz parte, suas transformações, os fenômenos que o modificam, os sujeitos e os seres que habitam o planeta, e as relações que se estabelecem entre os sujeitos.
Esse campo faz menção ainda às experiências da \textbf{criança} com o conhecimento da lógica matemática, não como disciplina, mas como uma linguagem que possibilite a compreensão do mundo em que vive, possibilitando assim uma aprendizagem significativa. > [... ] a Educação \textbf{Infantil} precisa promover experiências nas quais as \textbf{crianças} possam fazer observações, \textbf{manipular} objetos, investigar e explorar seu entorno, levantar hipóteses e consultar fontes de informação para buscar respostas às suas \textbf{curiosidades} e indagações.
Assim, a instituição escolar está criando oportunidades para que as \textbf{crianças} ampliem seus conhecimentos do mundo físico e sociocultural e possam utilizá -los em seu cotidiano ( BRASIL, 2017a, p. 41 ).
Na interligação com o Ensino Fundamental, o campo de experiência “ Espaços, tempos, quantidades, relações e transformações ” está essencialmente conectado com o eixo estruturante “ Espaço/Tempo e suas Transformações ”, pois é pela compreensão da relação do sujeito com o mundo ; da dinâmica das relações sociais que envolvem os aspectos sociais, políticos, culturais, afetivos, econômicos ; das mudanças histórico-sociais promovidas pela ação do homem sobre o tempo/espaço vivido que se constitui e se amplia o processo de aprendizagem. [ ^21 ] : Os eixos estruturantes do ensino fundamental podem ser estudados minuciosamente no item que discorre especificamente acerca deles.

% matched lemmas: adulto, balbucio, bebê, criança, cuidado, curiosidade, escuta, gesto, infantil, manipular, pequeno, pintura, rotina, som
\end{textsample}
